\section{Step 1: Extract asymptotic coefficients from a scattering problem}
\label{sec:textbook-scattering}
% BEGIN CURATED SUBJECT INDEX
\index{Jost function}
\index{Wronskian}
\index{wave packet}
\index{Lippmann--Schwinger equation}
\index{cross section}
\index{partial-wave expansion}
\index{boundary condition!incoming and outgoing}
\index{Zel'dovich regularization}
% END CURATED SUBJECT INDEX

\calculationroute
  {The Schr\"odinger equation, wave packets, and elementary partial-wave
  decomposition.}
  {Define the comparison dynamics, flux normalization, Jost bases, and
  incoming/outgoing coefficients that will enter the exact-WKB connection
  problem.}
  {Given a short-range potential, write the Lippmann--Schwinger and partial-
  wave problems with consistent sheet, flux, and Wronskian conventions.}

This section fixes the elementary scattering-theory conventions used in the
rest of the book.  The order of presentation follows the traditional route
from wave packets and cross sections to partial waves, integral equations,
scattering matrices, and Jost functions.  Useful complementary treatments are
given by Sasakawa, Taylor, Perelomov and Zel'dovich, and Rakityansky
\cite{Sasakawa:1991,Taylor:1972,PerelomovZeldovich:1998,
Rakityansky:2022Jost}.  These references record provenance and alternative
conventions; the definitions and derivations needed below are developed here.
The steps that begin as ordinary collision theory will later become analytic
continuation, resonance normalization, and complex spectral deformation.

\subsection{Two-body reduction and the class of potentials}
\label{subsec:two-body-reduction}
% BEGIN CURATED SUBJECT INDEX
\index{Jost function}
\index{wave operator}
\index{scattering!Coulomb}
\index{radial equation}
% END CURATED SUBJECT INDEX

Consider two spinless particles with masses $m_1$ and $m_2$ interacting
through a translationally invariant potential $V(\bm r_1-\bm r_2)$.  With
center-of-mass and relative coordinates
\begin{equation}
  \bm R=\frac{m_1\bm r_1+m_2\bm r_2}{m_1+m_2},
  \qquad
  \bm r=\bm r_1-\bm r_2,
  \label{eq:cm-relative-coordinates}
\end{equation}
the Hamiltonian separates as
\begin{equation}
  \vsub{H}{tot}
  =-\frac{\hbar^2}{2(m_1+m_2)}\nabla_{\bm R}^2+H,
  \qquad
  H=H_0+V,
  \qquad
  H_0=-\frac{\hbar^2}{2\mu}\nabla_{\bm r}^2,
  \label{eq:two-body-relative-H}
\end{equation}
where $\mu=m_1m_2/(m_1+m_2)$ is the reduced mass.  The center-of-mass
plane wave is dynamically trivial, so the collision problem is the spectral
problem for $H$.

We first assume a real, local, central potential $V(r)$ that is nonsingular
enough at the origin and short ranged at infinity.  A representative
sufficient condition for the simplest one-channel statements is
\begin{equation}
  \int_0^\infty \rmd r\,r\lvert V(r)\rvert<\infty .
  \label{eq:short-range-condition}
\end{equation}
The exact hypotheses vary with the result: existence of wave operators,
analyticity of Jost functions, and dilation analyticity require different
weights or regularity.  Coulomb forces violate
Eq.~\eqref{eq:short-range-condition}; one must then replace free asymptotic
waves by Coulomb waves and factor out the known Coulomb phase.  Nonlocal,
coupled-channel, optical, and energy-dependent interactions are introduced
only after the one-channel structure is clear.

For a central potential, separation in spherical harmonics gives
\begin{equation}
  \psi(\bm r)=\sum_{lm}\frac{u_l(r)}{r}Y_{lm}(\widehat{\bm r}),
  \label{eq:partial-wave-decomposition}
\end{equation}
and the reduced radial equation
\begin{equation}
  \left[-\frac{\hbar^2}{2\mu}\frac{\rmd^2}{\rmd r^2}
  +\frac{\hbar^2l(l+1)}{2\mu r^2}+V(r)\right]u_l(r)
  =E u_l(r) .
  \label{eq:radial-schrodinger-textbook}
\end{equation}
Regularity of the full wave function requires
$u_l(r)=O(r^{l+1})$ at the origin.  The second boundary condition at infinity
distinguishes bound, scattering, virtual, and resonant solutions.  This is the
first indication that these are not four unrelated differential equations,
but four boundary-value realizations of the same radial equation.

\subsection{Wave packets, asymptotic states, and wave operators}
\label{subsec:wave-packets-operators}
% BEGIN CURATED SUBJECT INDEX
\index{locally convex space}
\index{wave packet}
\index{wave operator}
\index{scattering matrix}
\index{boundary condition!incoming and outgoing}
\index{Gamow--Siegert state}
% END CURATED SUBJECT INDEX

A plane wave is not normalizable and does not by itself describe a collision
prepared at a finite time.  Begin instead with a free wave packet
\begin{equation}
  \phi(\bm r,t)
  =\int\frac{\rmd^3k}{(2\pi)^3}\,
  a(\bm k)\rme^{\rmi\bm k\cdot\bm r-\rmi E_k t/\hbar},
  \qquad
  E_k=\frac{\hbar^2k^2}{2\mu},
  \label{eq:free-wave-packet}
\end{equation}
with $a(\bm k)$ concentrated near an incident momentum $\bm k_0$.  The
interacting packet is required to approach a free packet before or after the
collision.  When the strong limits exist, this statement is encoded by the
M{\o}ller wave operators
\begin{equation}
  \Omega^{(\pm)}
  =\operatorname*{s-lim}_{t\to\mp\infty}
  \rme^{\rmi Ht/\hbar}\rme^{-\rmi H_0t/\hbar}\vsub{P}{ac}(H_0) .
  \label{eq:moller-wave-operators}
\end{equation}
Here $\vsub{P}{ac}(H_0)$ projects onto the absolutely continuous
subspace.  The state $\Omega^{(+)}\phi$ has the prescribed incoming free
packet in the remote past, whereas $\Omega^{(-)}\phi$ has the prescribed
outgoing free packet in the remote future.  The scattering operator is
\begin{equation}
  \Scattering=\Omega^{(-)\dagger}\Omega^{(+)} .
  \label{eq:scattering-operator-wave}
\end{equation}
If the wave operators are isometric and asymptotically complete,
$\Scattering$ is unitary on the free scattering space.  Asymptotic
completeness means that all states orthogonal to bound states eventually
resolve into scattering channels.  It is a theorem under specified
hypotheses, not an automatic consequence of writing $H=H_0+V$.

The stationary scattering states used below are distributional limits of
these packets.  Their delta-function normalization is therefore a shorthand
for a wave-packet construction.  This distinction matters for resonances:
neither an exact plane wave nor a Gamow state is an ordinary Hilbert-space
vector, but both define continuous functionals on a suitable test space.

\subsection{Free resolvent and the Lippmann--Schwinger boundary condition}
\label{subsec:free-resolvent-ls}
% BEGIN CURATED SUBJECT INDEX
\index{resolvent}
\index{Lippmann--Schwinger equation}
\index{scattering!on-shell}
\index{boundary condition!incoming and outgoing}
% END CURATED SUBJECT INDEX

Let
\begin{equation}
  \Resolvent_0(z)=(z-H_0)^{-1} .
  \label{eq:free-resolvent-textbook}
\end{equation}
For $E=\hbar^2k^2/(2\mu)>0$, the coordinate kernel of its outgoing boundary
value is
\begin{equation}
  \bra{\bm r}\Resolvent_0(E+\rmi0)\ket{\bm r'}
  =-\frac{\mu}{2\pi\hbar^2}
  \frac{\rme^{\rmi k\lvert\bm r-\bm r'\rvert}}
  {\lvert\bm r-\bm r'\rvert} .
  \label{eq:free-green-outgoing}
\end{equation}
The sign $+\rmi0$ is not a formal decoration.  Fourier inversion places the
radial momentum pole on the side of the contour that produces an outgoing
spherical wave.  The opposite boundary value produces an incoming wave.

For an incident plane wave $\braket{\bm r\vert\bm k}=\rme^{\rmi\bm
k\cdot\bm r}$, the outgoing Lippmann--Schwinger equation is
\begin{equation}
  \ket{\psi_{\bm k}^{(+)}}
  =\ket{\bm k}+\Resolvent_0(E_k+\rmi0)V
  \ket{\psi_{\bm k}^{(+)}} .
  \label{eq:ls-coordinate-foundation}
\end{equation}
Substitution of Eq.~\eqref{eq:free-green-outgoing} and expansion at large
$r$ give
\begin{equation}
  \psi_{\bm k}^{(+)}(\bm r)
  \sim \rme^{\rmi\bm k\cdot\bm r}
  +f(\widehat{\bm r},\bm k)\frac{\rme^{\rmi kr}}{r},
  \qquad r\to\infty,
  \label{eq:scattering-asymptotic-3d}
\end{equation}
where
\begin{equation}
  f(\widehat{\bm r},\bm k)
  =-\frac{\mu}{2\pi\hbar^2}
  \int\rmd^3r'\,
  \rme^{-\rmi k\widehat{\bm r}\cdot\bm r'}
  V(\bm r')\psi_{\bm k}^{(+)}(\bm r') .
  \label{eq:scattering-amplitude-integral}
\end{equation}
Thus the radiation condition is already contained in the resolvent boundary
value.  Later, analytic continuation of precisely this boundary value turns
the outgoing solution into a resonance state.

The transition operator
\begin{equation}
  T(z)=V+V\Resolvent_0(z)T(z)
  \label{eq:T-foundation}
\end{equation}
packages Eq.~\eqref{eq:scattering-amplitude-integral} as
\begin{equation}
  f(\widehat{\bm r},\bm k)
  =-\frac{\mu}{2\pi\hbar^2}
  \bra{k\widehat{\bm r}}T(E_k+\rmi0)\ket{\bm k} .
  \label{eq:f-T-relation}
\end{equation}
Normalization-dependent factors change if momentum states are normalized
differently, but pole positions and on-shell observables do not.

\subsection{Flux, cross sections, and the optical theorem}
\label{subsec:flux-cross-sections}
% BEGIN CURATED SUBJECT INDEX
\index{resolvent}
\index{optical theorem}
\index{cross section}
\index{boundary condition!incoming and outgoing}
\index{Feshbach projection}
\index{effective Hamiltonian}
\index{optical potential}
% END CURATED SUBJECT INDEX

For a wave function obeying a real-potential Schrodinger equation, the
probability current is
\begin{equation}
  \bm j[\psi]
  =\frac{\hbar}{2\mu\rmi}
  \left(\psi^*\nabla\psi-\psi\nabla\psi^*\right) .
  \label{eq:probability-current-3d}
\end{equation}
The incident plane wave in Eq.~\eqref{eq:scattering-asymptotic-3d} carries
flux $\hbar\bm k/\mu$, and the outgoing spherical wave carries radial flux
proportional to $\lvert f\rvert^2/r^2$.  Their ratio gives
\begin{equation}
  \frac{\rmd\sigma}{\rmd\Omega}
  =\lvert f(\widehat{\bm r},\bm k)\rvert^2 .
  \label{eq:differential-cross-section}
\end{equation}
Conservation of total flux and interference between the incident wave and the
forward part of the scattered wave imply the optical theorem
\begin{equation}
  \vsub{\sigma}{tot}
  =\frac{4\pi}{k}\im f(\widehat{\bm k},\bm k) .
  \label{eq:optical-theorem}
\end{equation}
In operator form it follows from the discontinuity of the free resolvent and
the identity
\begin{equation}
  T(E+\rmi0)-T(E-\rmi0)
  =-2\pi\rmi\,T(E-\rmi0)\delta(E-H_0)T(E+\rmi0) .
  \label{eq:T-unitarity-identity}
\end{equation}
Equation~\eqref{eq:T-unitarity-identity} is the precise reason that the
imaginary part of an effective Hamiltonian represents lost flux into open
channels while the full theory remains unitary.

If $V$ is replaced by a complex optical potential, elastic flux alone is not
conserved.  The deficit represents channels removed from the reduced model.
One then has $\lvert S_l\rvert\leq1$ rather than
$\lvert S_l\rvert=1$, and the reaction cross section measures the missing
flux.  This elementary example anticipates the Feshbach effective
Hamiltonian of Section~\ref{sec:scattering}.

\subsection{Partial waves, phase shifts, and threshold data}
\label{subsec:partial-waves-phase-shifts}
% BEGIN CURATED SUBJECT INDEX
\index{cross section}
\index{partial-wave expansion}
\index{phase shift}
\index{threshold law}
\index{boundary condition!incoming and outgoing}
% END CURATED SUBJECT INDEX

The plane-wave expansion
\begin{equation}
  \rme^{\rmi\bm k\cdot\bm r}
  =4\pi\sum_{lm}\rmi^l j_l(kr)
  Y_{lm}(\widehat{\bm r})Y_{lm}^*(\widehat{\bm k})
  \label{eq:rayleigh-expansion}
\end{equation}
reduces central-potential scattering to independent radial channels.  For a
real short-range potential, the large-$r$ form of the regular radial solution
can be normalized as
\begin{equation}
  u_l(k,r)\sim
  \frac{1}{2\rmi}
  \left[S_l(k)\rme^{\rmi(kr-l\pi/2)}
  -\rme^{-\rmi(kr-l\pi/2)}\right] .
  \label{eq:partial-wave-asymptotic}
\end{equation}
Elastic unitarity gives
\begin{equation}
  S_l(k)=\rme^{2\rmi\delta_l(k)},
  \qquad \delta_l(k)\in\real,
  \label{eq:S-phase-shift}
\end{equation}
and the scattering amplitude becomes
\begin{equation}
  f(\vartheta)
  =\frac{1}{2\rmi k}\sum_{l=0}^{\infty}(2l+1)
  \left[S_l(k)-1\right]P_l(\cos\vartheta) .
  \label{eq:partial-wave-amplitude}
\end{equation}
Consequently
\begin{align}
  \vsub{\sigma}{el}
  &=\frac{4\pi}{k^2}\sum_{l=0}^{\infty}(2l+1)
  \sin^2\delta_l,
  \label{eq:partial-wave-cross-section}\\
  f_l(k)&=\frac{S_l(k)-1}{2\rmi k}
  =\frac{1}{k\cot\delta_l(k)-\rmi k}
  \qquad (l=0) .
  \label{eq:s-wave-amplitude}
\end{align}

Near a short-range threshold, the effective-range expansion for the $s$ wave
is
\begin{equation}
  k\cot\delta_0(k)
  =-\frac{1}{a_0}+\frac{1}{2}r_0k^2+O(k^4),
  \label{eq:effective-range-expansion}
\end{equation}
where $a_0$ and $r_0$ are the scattering length and effective range.  A large
$\lvert a_0\rvert$ signals a pole close to threshold, but the pole can be a
shallow bound state or a virtual state depending on its sheet.  This is why a
large cross section alone does not classify the singularity.

The standing-wave or reactance matrix is particularly useful on the real
axis.  In one partial wave set $K_l=\tan\delta_l$; then
\begin{equation}
  S_l=\frac{1+\rmi K_l}{1-\rmi K_l} .
  \label{eq:S-K-relation}
\end{equation}
For a real potential $K_l$ is real, while $S_l$ is unitary.  The $T$, $K$,
and $S$ matrices therefore emphasize, respectively, outgoing boundary
values, principal-value standing waves, and the mapping from incoming to
outgoing channel amplitudes.  Their singularities are related, but a pole of
$K_l$ on the real axis is not by itself the invariant definition of a
resonance.

\subsection{Jost solutions as a common coordinate basis}
\label{subsec:jost-foundation}
% BEGIN CURATED SUBJECT INDEX
\index{Jost function}
\index{Wronskian}
\index{scattering!multichannel}
\index{boundary condition!incoming and outgoing}
\index{scattering!Coulomb}
% END CURATED SUBJECT INDEX

The regular solution used in Eq.~\eqref{eq:partial-wave-asymptotic} builds the
origin boundary condition into the differential equation.  Jost solutions
instead build the asymptotic boundary conditions into it.  Let
$f_l(k,r)$ and $f_l(-k,r)$ be solutions of
Eq.~\eqref{eq:radial-schrodinger-textbook} satisfying
\begin{equation}
  f_l(\pm k,r)\sim\rme^{\pm\rmi(kr-l\pi/2)},
  \qquad r\to\infty,
  \label{eq:jost-foundation-asymptotic}
\end{equation}
where an $l$-dependent constant phase has been absorbed into the definition.
For $k\ne0$ they form a basis of the two-dimensional solution space.  If
$\varphi_l(k,r)$ is regular at the origin, it has the expansion
\begin{equation}
  \varphi_l(k,r)
  =\frac{1}{2\rmi k}
  \left[F_l(-k)f_l(k,r)-F_l(k)f_l(-k,r)\right] .
  \label{eq:jost-basis-foundation}
\end{equation}
With the Wronskian convention
$W[u,v]=uv'-u'v$, the coefficients are
\begin{equation}
  F_l(\pm k)=W[f_l(\pm k,\cdot),\varphi_l(k,\cdot)] .
  \label{eq:jost-wronskian-definition}
\end{equation}
The normalization in Eqs.~\eqref{eq:jost-basis-foundation} and
\eqref{eq:jost-wronskian-definition} is consistent with
\begin{equation*}
  W[f_l(k),f_l(-k)]=-2\rmi k.
\end{equation*}

The Jost function is thus not an auxiliary spectral determinant introduced
after solving the problem.  It is the coefficient of the incoming Jost basis
vector in the solution regular at the origin.  The quotient
\begin{equation}
  S_l(k)=\frac{F_l(-k)}{F_l(k)}
  \label{eq:S-Jost-foundation}
\end{equation}
follows immediately.  A zero of $F_l(k)$ removes the incoming component,
leaving a solution that is regular at the origin and outgoing at infinity.
Depending on the position of the zero in the $k$ plane, the same algebraic
condition describes a bound, virtual, or resonant state.  This is the central
unification emphasized in the Jost-function approach
\cite{Rakityansky:2022Jost}.

For a finite-range potential the Jost solutions and Jost function possess a
large analytic domain in $k$; for exponentially decreasing potentials the
domain is restricted by the decay rate.  The point $k=0$ requires separate
threshold analysis, and Coulomb interactions require Coulomb-modified Jost
functions.  In an $N$-channel problem, $F_l(k)$ becomes a Jost matrix and
bound or resonant states satisfy
\begin{equation}
  \det F(E)=0 .
  \label{eq:multichannel-jost-determinant}
\end{equation}
Each channel momentum $k_a(E)$ supplies its own square root, so the determinant
lives on a multi-sheeted energy surface.  Keeping these sheets explicit is
essential when a resonance crosses a threshold or coalesces with another
pole.
